\section{Artifact status, reproducibility, and integration checklist}
\label{app:status}

\subsection{Delivered files}

\begin{longtable}{@{}p{0.42\textwidth}p{0.51\textwidth}@{}}
\toprule
Path & Purpose \\
\midrule
\endfirsthead
\toprule
Path & Purpose \\
\midrule
\endhead
\path{paper/scr_ipra_framework.tex} & complete LaTeX driver \\
\path{paper/sections/*.tex} & editable report sections \\
\path{paper/scr_ipra_framework.pdf} & compiled report \\
\path{lean/ScopedCertificate/Kernel.lean} & no-placeholder structural kernel candidate \\
\path{lean/ScopedCertificate/Blueprint.lean} & host metadata and commented object targets \\
\path{lean/FORMALIZATION_PLAN.md} & dependency-ordered implementation and acceptance gates \\
\path{reference/scr_reference.py} & exact finite encodings and examples \\
\path{reference/test_scr_reference.py} & eight standard-library regression tests \\
\path{examples/demo.json} & generated bisection, path, observer, and height-list evidence \\
\path{schema/scr_program.schema.json} & machine-readable program metadata schema \\
\path{metadata/THEOREM_LEDGER.md} & claim-by-claim status \\
\path{metadata/DIFFICULTY_MATRIX.csv} & obstacle/resolution mapping \\
\path{metadata/SOURCE_AUDIT.md} & pinned repository sources \\
\path{validation/} & test, PDF, rendering, and static-audit records \\
\path{MANIFEST.sha256} & checksums of delivered files \\
\bottomrule
\end{longtable}

\subsection{Reproduction commands}

From the archive root:
\begin{verbatim}
python3 -m unittest -v reference/test_scr_reference.py
python3 reference/scr_reference.py > examples/demo.json
bash build.sh
\end{verbatim}
The build script runs the Python tests and compiles the PDF.  Lean verification
requires copying \code{lean/} into a checkout of the pinned \repo\ tree and
following \path{lean/FORMALIZATION_PLAN.md}.

\subsection{Integration checklist}

A packet is ready for repository integration under SCR when all applicable
boxes are checked.

\begin{itemize}[label=$\square$]
\item The public theorem is stated as a literal \(\Prf\;\PRAAx\) judgment.
\item Every host algorithm has a fixed PR symbol and object defining equations.
\item Every bounded logical condition has a reflected checker and proof.
\item The qf proof for each induction invariant is explicit.
\item Each certificate is marked sealed, observed, or named.
\item Every named certificate has a direct producer or explicit code cutoff.
\item Whole paths are used where greedy local selection can fail.
\item Noncanonical outputs have an observer-invariance or diameter theorem.
\item Exact labels include a uniqueness or separation premise.
\item Cardinality bounds are not used as code cutoffs.
\item Missing moduli, heights, covers, and constants occur as boundary formulas.
\item Reusable lemmas work under local contexts and at arbitrary terms.
\item New declarations pass focused builds and the repository axiom audit.
\item Executable evidence is not described as an object proof.
\item The source pin and all remaining assumptions are recorded.
\end{itemize}

\subsection{Claim policy}

The generic human proofs in this document are intended to be complete, but
formal scoping details can still reveal errors during Lean elaboration.  The
archive therefore distinguishes mathematical confidence from kernel status.
If a candidate theorem fails to elaborate, the report's claim remains a human
proof proposal until repaired and checked.  Conversely, a passing Python test
confirms only the finite reference algorithm, not its universal object theorem.

This status discipline is part of the framework itself.  A proof-producing
architecture is trustworthy only when its own artifacts obey the same
producer/verifier separation that it recommends for mathematics.
